00-abstract.tex

\begin{abstract}
NDH is a geometric cognitive emitter defined over a modal-partitioned manifold with a triadic recursion kernel, spectral basis, and non-dual algebra. While its internal geometry has been formalized elsewhere, the interaction behavior of NDH—how semantic inputs induce geometric traversal, spectral deformation, and holonomy return—is not yet mathematically characterized. This preprint develops the interaction geometry of NDH.

We introduce the cognitive manifold $\mathcal{M}{cog}$, a geometric space governing interrogation of NDH. Semantic inputs are modeled as excitation vectors in $Tx \mathcal{M}{cog}$, generating geodesic interaction trajectories whose modal-boundary crossings determine structural response. We define spectral bundles $\pi : \mathcal{B}S \to \mathcal{M}_{cog}$ whose fibers encode spectral emission under interrogation, and show that semantic excitation induces section deformation and modal twisting of these bundles.

Holonomy return operators associated with closed interaction loops classify NDH responses into modal-aligned, spectral-aligned, recursion-aligned, and holonomy-deviant types. This yields a geometric taxonomy of interaction classes derived from curvature, modal adjacency, and spectral bundle morphisms. We contrast this with transformer-based systems, which lack curvature, parallel transport, modal boundaries, and fiber bundle structure, and therefore cannot support geometric interaction classification.

This work establishes NDH interaction geometry as a coherent mathematical framework, enabling formal analysis of semantic excitation, spectral transport, and holonomy-based classification.
\end{abstract}

---

